\documentclass[a4paper,12pt,twoside]{scrreprt}
% Autor der Vorlage: Klaus Rheinberger, FH Vorarlberg, 2017-02-20

% Pakete:
\usepackage[utf8]{inputenc}
\usepackage[T1]{fontenc} % Silbentrennung bei Sonderzeichen
\usepackage{graphicx} % Bilder einbinden
\usepackage{wrapfig} % Bilder positionieren
\usepackage[english]{babel} 
\usepackage{minted} % Code Highlighting/Import
\usepackage{csquotes} % Anführungszeichen und Zitieren
\usepackage[bindingoffset=8mm]{geometry} % Bindeverlust von 8mm einbeziehen
\usepackage{caption} % Abbildungslegenden
\usepackage{setspace} % Zeilenabstand
\usepackage[style=alphabetic,citestyle=alphabetic,backend=biber]{biblatex} % Literaturverweise
\usepackage[
    linktocpage=true,
    pdfauthor={Dominik Böckle & Dominic Luidold},
    pdftitle={Translating the concepts of Haskell into the object-oriented programming language Java}
]{hyperref} % Links
\usepackage{acronym} 
\usepackage{eurosym}
\usepackage{float}

% Einstellungen:
\captionsetup{format=hang, justification=raggedright}
\addbibresource{Sources.bib}
\setcounter{secnumdepth}{4}
\renewcommand{\listingscaption}{Source Code}
\renewcommand\listoflistingscaption{List of Source Codes}

% Dokumentenbeginn
\begin{document}
\onehalfspacing % Zeilenabstand 1,5

% Titelblatt:
% \newpage\mbox{}\newpage
\pagenumbering{Roman}
\cleardoublepage % force output to a right page
\thispagestyle{empty}
\begin{titlepage}
    \begin{flushright}
    \includegraphics[width=0.4\linewidth]{Logo-A3.jpg}
    \end{flushright}
    \begin{flushleft}
    \vspace{0.5cm}
    \section*{Translating the concepts of Haskell into the object-oriented programming language Java}
    \subsection*{Written Elaboration}

    \vspace{4cm}
    Vorarlberg University of Applied Sciences\newline
    Computer Science MSc

    \vspace{0.5cm}

    Supervised by\newline
    Jonathan Thaler, Ph.D

    \vspace{0.5cm}

    Submitted by\newline
    Dominik Böckle,\newline
    Dominic Luidold\newline\newline
    Dornbirn, January 2022
    \end{flushleft}
\end{titlepage}

% Inhaltsverzeichnis:
\cleardoublepage % force output to a right page
\setcounter{tocdepth}{2}
\pdfbookmark{\contentsname}{toc}
\tableofcontents

% Abkürzungsverzeichnis:
\clearpage
\phantomsection
\addcontentsline{toc}{chapter}{List of Abbreviations}
\chapter*{List of Abbreviations}
\begin{acronym}
  \acro{ADT}{Algebraic Data Types}
  \acro{OOP}{Object-oriented programming}
\end{acronym}

% Quellcodeverzeichnis
\clearpage
\phantomsection
\addcontentsline{toc}{chapter}{List of Source Codes}
\listoflistings

\cleardoublepage
\pagenumbering{arabic}
\chapter{Introduction}
This written elaboration contains various concepts of functional programming that are covered in the lecture \enquote{Concepts of Higher Programming Languages} in the first semester Master of Science at the Vorarlberg University of Applied Sciences. The concepts and patterns covered in this course are implemented in this paper using the \ac{OOP} language Java. Each implementation is complemented by a short explanation.

In addition, a \texttt{StateMonad} is implemented in Java and an example of how to use the Monad implemnetation is provided.

\chapter{Implementation of the various concepts}
This chapter covers the implementation of eight functional programming concepts and their respective \acs{OOP} counterparts written in Java. The concepts and patterns are each accompanied by a brief explanation as well as a working Java example.

\section{Immutable Data}
\textit{Immutable Data} is a principle that states that once data has been created, it subsequently cannot be changed. In general, data and objects should only be mutable if there is a valid reason for doing so.

\begin{listing}[ht]
    \inputminted[fontsize=\footnotesize,linenos]{java}{./code/ImmutableIntegerPair.java}
    \caption[Example of an immutable data structure]{Example of an immutable data structure}
    \label{code:immutable}
\end{listing}

\textit{Source Code \ref{code:immutable}} on page \pageref{code:immutable} shows how an immutable data structure may look. The variables \texttt{x} and \texttt{y}, which contain the data, can no longer be modified after the value has initially been set as a result of the \texttt{final} keyword. Additionally, the variables cannot be changed from outside the class as there are no setters present.

There is, however, still a possibility to make the data structure mutable: Extending the \texttt{ImmutableIntegerPair} allows one to extend the functionality of the data structure and store mutable data. To prevent this, the keyword \texttt{final} can be used again, this time at class level as opposed to variable level.

\medskip

Furthermore, there are also ways to make lists immutable when the list data type is the preferred/required way of storing data. Java itself provides a static method \texttt{Collections::unmodifiableList()} which prevents the data from being modified. However, the data type that is stored in the list must also be immutable to achieve immutability.

\section{Type Variables}
As far as the concept of \textit{Type Variables} is concerned, Java - being a strongly typed programming language - enables the typing of variables and the enforcement of a certain type throughout the programme flow in various ways:

\begin{itemize}
    \item Each variable has to have an accompanying data type when getting declared, e.g. \texttt{int x = 1;}
    \item Java natively supports Generics (as can be seen in \textit{Source Code \ref{code:ADT}} on page \pageref{code:ADT}) which allows types to be enforced based on the specific use case at compile time. Constructs such as \texttt{<? extends Integer>} additionally allow the restriction of which types are \enquote{allowed}.
    \item Checking types at run time is possible using \texttt{instanceof} when there are several types possible at one point in a time, for example when having a common interface and different classes inheriting that interface.
\end{itemize}

\clearpage

\section{Higher-Order Functions}
A \textit{Higher-Order Function} is a function that either takes a function as an argument or has a function as a return value. Implementing this concept has been possible since Java 8 - an example for that is shown in \textit{Sourcode \ref{code:higherordered}} on page \pageref{code:higherordered}. The method \texttt{camelize} capitalises the first character of a string and returns the required method/function calls when called. This method also makes use of \textit{Lambdas}, which is explained in more detail in a subsequent chapter.

\begin{listing}[ht]
    \inputminted[fontsize=\footnotesize,linenos,breaklines,breakanywhere]{java}{./code/HigherOrderFunctions.java}
    \caption[Example of a Higher-Order Function in Java 8]{Example of a Higher-Order Function in Java 8}
    \label{code:higherordered}
\end{listing}

\clearpage

\section{Lambda Expressions}
The method \texttt{listLambdas} in \textit{Sourcode \ref{code:Lambda}} on page \ref{code:Lambda} increments the \texttt{Integer} values stored in a \texttt{List} by one. This is achieved by the rather short code fragment \texttt{map(x -> ++x)}.

The same functionality can be achieved be creating an anonymous \texttt{Function} instance within the \texttt{map} method (not shown in this example), which has been the recommended and only way until to Java 8 to implement such functionality. However, Lambda Expressions have the advantage of code that is usually more readable and requires less boilerplate code.

\begin{listing}[ht]
    \inputminted[fontsize=\footnotesize,linenos]{java}{./code/LambdaExpressions.java}
    \caption[Example of a Lambda Expression]{Example of a Lambda Expression}
    \label{code:Lambda}
\end{listing}

\clearpage

\section{Currying}
\textit{Currying} is the conversion of a function with multiple arguments into a sequence of functions with one argument each. \textit{Source Code \ref{code:Currying}} on page \pageref{code:Currying} is inspired by \cite{Robertson_currying_2018} in which multiplication was performed by currying. This is emulated with the function \texttt{multCurry} where the multiplications can be concatenated.

\begin{listing}[ht]
    \inputminted[fontsize=\footnotesize,linenos,breaklines,breakanywhere]{java}{./code/Currying.java}
    \caption[Example of Currying]{Example of Currying}
    \label{code:Currying}
\end{listing}

\clearpage

\section{Function Composition and Streaming}

\subsection{Function Composition}
\textit{Function Composition} is applying the pattern of combining multiple separate functions into a single function. The output of the first function is then used as input for the second function and so on. To be able to demonstrate this in an example, \textit{Source Code \ref{code:Compose}} on page \pageref{code:Compose} containts the functions \texttt{doubleInt} and \texttt{substractOne} - both based on the use of Lambda Expressions. 

An important part of Function Composition is that order matters. While using  \texttt{compose(<fn>)} results in the supplied function (\texttt{substractOne}) being executed before the initial function (\texttt{doubleInt}), using \texttt{andThen(<fn>)} has the opposite result.

\begin{listing}[ht]
    \inputminted[fontsize=\footnotesize,linenos,breaklines,breakanywhere]{java}{./code/FunctionComposition.java}
    \caption[Example of Function Composition and that order matters]{Example of Function Composition and that order matters}
    \label{code:Compose}
\end{listing}

\clearpage

\subsection{Streaming}
A \textit{Stream} or \textit{Streaming} represents a sequence of objects that are accessed in sequential order. One of the best known and most frequently used Streams in Java is probably the \texttt{FileStream}, which is used to read files and file contents. \textit{Source Code \ref{code:Streaming}} on page \pageref{code:Streaming} shows this in a small, but only theoretical example.

\begin{listing}[ht]
    \inputminted[fontsize=\footnotesize,linenos,breaklines]{java}{./code/Streaming.java}
    \caption[Example of Streaming]{Example of Streaming.}
    \label{code:Streaming}
\end{listing}

\clearpage

\section{Algebraic Data Types}
\textit{\ac{ADT}} consist of several variants or \enquote{flavors}, similar to a Java Enum, but these different \enquote{flavors} can each have different properties and/or methods. To quickly summarise what we expect from an \ac{ADT}:
\begin{itemize}
    \item The expressions have different amounts and types of properties and methods.
    \item At compile time, it is checked whether all have been taken into account.
\end{itemize}

Unfortunately, Java does not support \ac{ADT} natively, so we had to implement this ourselves. For this, as described in \cite{MAINIERO_algebraic_2020}, the \textit{Visitor} and \textit{Sealed Class} design patterns are helpful. Since Java 15, the \texttt{sealed} keyword is available, but the following example nevertheless includes a manually implemented \textit{sealed} functionality. \textit{Source Code \ref{code:ADT}} on page \pageref{code:ADT} and before shows a corresponding Java implementation of the \textit{Dead or Alive} State from the lecture notes\footnote{\href{https://homepages.fhv.at/thjo/lecturenotes/concepts/declaring-types.html\#algebraic-data-types-1}{Lecture Notes, Jonathan Thaler, Ph.D (visited on 2022/01/06)}}

\inputminted[fontsize=\footnotesize,linenos,breaklines,breakanywhere]{java}{./code/AlgebraicDataTypes.java}

\begin{center}
    \captionof{listing}{\enquote{Dead or Alive} example demonstrating \acs{ADT}\label{code:ADT}}
\end{center}

\section{Pure and Impure Side Effects}
\textit{Pure Side Effects} are effects which are solely caused by a method or calculation. They are always deterministic, as they do not communicate with the \enquote{outside world}. \textit{Source Code \ref{code:Effects}} on page \pageref{code:Effects} shows the change of the \texttt{State} within the method \texttt{pure}.

\textit{Impure Side Effects} are all effects caused by interacting with the \enquote{outside world}. This can be seen in \textit{Source Code \ref{code:Effects}} within the \texttt{impure} method which awaits user action/input.

\begin{listing}[ht]
    \inputminted[fontsize=\footnotesize,linenos,breaklines,breakanywhere]{java}{./code/SideEffects.java}
    \caption[Example of Side Effects in Java]{Example of Side Effects in Java}
    \label{code:Effects}
\end{listing}

\clearpage

\section{State Monad in Java}
The following section contains an implementation of the \textit{State Monad} in Java. During the implementation of the Monad, the work of \cite{muzietto_state_2014} was of great support. A question that might arise: \enquote{What is a Monad?}. Matt Fowler describes it as follows: 
\begin{quote}
    Think of monads as an object that wraps a value and allows us to apply a set of transformations on that value and get it back out with all the transformations applied. \cite{fowler_understanding_2015}
\end{quote}

\subsection{Implementing the State Monad}
To develop a working State Monad in Java, both \textit{lift} and \textit{bind} had to be implemented as well as a \texttt{StateTuple} that stores the state as well as the values as content. \textit{Source Code \ref{code:StateMonad}} on page \pageref{code:StateMonad} shows the class \texttt{StateMonad<S, C>} with an implementation of \textit{bind}.

\begin{listing}[ht]
    \inputminted[fontsize=\footnotesize,linenos,breaklines,breakanywhere]{java}{./code/statemonad/StateMonad.java}
    \caption[State Monad implementation]{State Monad implementation}
    \label{code:StateMonad}
\end{listing}

\clearpage

\textit{Source Code \ref{code:StateTuple}} on page \pageref{code:StateTuple} shows the implementation representing the state and the content of the data within the State Monad.

\begin{listing}[ht]
    \inputminted[fontsize=\footnotesize,linenos,breaklines,breakanywhere]{java}{./code/statemonad/StateTuple.java}
    \caption[\texttt{StateTuple} implementation]{\texttt{StateTuple} implementation}
    \label{code:StateTuple}
\end{listing}
\clearpage

\subsection{Implementing Tree Labelling}
\textit{Source Code \ref{code:TreeLabelling}} on page \pageref{code:TreeLabelling} and before shows the implementation of a \textit{Binary Tree} labelling, while \textit{Source Code \ref{code:Tree}} on page \pageref{code:Tree} shows the abstract class  \texttt{Tree} that defines the basic behaviour of the Binary Tree.

\inputminted[fontsize=\footnotesize,linenos,breaklines,breakanywhere]{java}{./code/statemonad/Labelling.java}
\captionof{listing}{Tree labelling implementation \label{code:TreeLabelling}}

\begin{listing}[ht]
    \inputminted[fontsize=\footnotesize,linenos,breaklines,breakanywhere]{java}{./code/statemonad/Tree.java}
    \caption[\texttt{Tree} implementation]{\texttt{Tree} implementation}
    \label{code:Tree}
\end{listing}

% Literaturverzeichnis:
\clearpage
\phantomsection
\addcontentsline{toc}{chapter}{Bibliography}
\printbibliography

\end{document}
